\lyseis{Β3.12}{Κύκλος και στοιχεία του}
\sol{1} α) $360$. β) $180$. γ) $90$. δ) $\delta=2\cdot7=14$ cm. ε) $\rho=18:2=9$ cm.
\sol{2} α) Σωστό. β) Λάθος. γ) Σωστό. δ) Σωστό. ε) Σωστό. στ) Λάθος, η διάμετρος είναι η μεγαλύτερη χορδή.
\sol{3} Ελεύθερη κατασκευή. Θα βάψεις το εσωτερικό του κύκλου ακτίνας $2$ cm με το πρώτο χρώμα, τον ίδιο τον κύκλο με το δεύτερο και ό,τι είναι έξω με το τρίτο.
\sol{4} Ελεύθερη κατασκευή. Η διάμετρος $\mathrm{ΑΑ'}$ περνά από το κέντρο και είναι $6$ cm, και η επίκεντρη γωνία $\widehat{\mathrm{ΑΟΒ}}$ έχει κορυφή το κέντρο $\mathrm{Ο}$.
\sol{5} Η ακτίνα του δεύτερου είναι $6:2=3$ cm. Οι κύκλοι απέχουν $3-2=1$ cm.
\sol{6} $\delta=13$ cm, μεγαλύτερη χορδή η διάμετρος, $13$ cm. Χορδή $14$ cm δεν υπάρχει, γιατί $14>13$.
\sol{7} Το ένα τόξο είναι $125^\circ$ και το άλλο $360^\circ-125^\circ=235^\circ$.
\sol{8} Ο Ηλίας ($12<15$) και η Ρήνη ($15=15$, στο όριο) έχουν σήμα. Ο Στέλιος ($18>15$) δεν έχει. Η μέγιστη απόσταση είναι η διάμετρος, $2\cdot15=30$ m.
\sol{9} Ο κύκλος των $360^\circ$ χωρίζεται σε $12$ ίσα τόξα: $360^\circ:12=30^\circ$. Ανάμεσα στο $12$ και στο $3$: $3\cdot30^\circ=90^\circ$. Από το $12$ ως το $5$: $5\cdot30^\circ=150^\circ$. Σε $60$ λεπτά κάνει $360^\circ$, άρα σε $10$ λεπτά $360^\circ:6=60^\circ$.
\sol{10} $360^\circ:8=45^\circ$. Τρία κομμάτια: $3\cdot45^\circ=135^\circ$.
\sol{11} Η διάμετρος, $2\cdot2{,}4=4{,}8$ m. Σχοινί $3$ m ναι ($3<4{,}8$). Σχοινί $5$ m όχι ($5>4{,}8$).
\sol{12} Το βελάκι είναι έξω από τον κύκλο των $10$ cm ($12>10$) και μέσα στον κύκλο των $15$ cm ($12<15$), δηλαδή ανάμεσα στον δεύτερο και στον τρίτο κύκλο.
\sol{13} α) Σχεδιάζουμε $\mathrm{ΑΒ}=5$ cm, κύκλους $(\mathrm{Α},4)$ και $(\mathrm{Β},3)$ και παίρνουμε ένα σημείο τομής για κορυφή. β) Όχι. Οι κύκλοι $(\mathrm{Α},4)$ και $(\mathrm{Β},2)$ με $\mathrm{ΑΒ}=7$ δεν τέμνονται, γιατί $4+2=6<7$.
\sol{14} Σχεδιάζουμε $\mathrm{ΑΒ}=5$ cm και δύο κύκλους ακτίνας $5$ cm, $(\mathrm{Α},5)$ και $(\mathrm{Β},5)$. Ένα σημείο τομής τους $\mathrm{Γ}$ έχει $\mathrm{ΑΓ}=5$ (ακτίνα του πρώτου) και $\mathrm{ΒΓ}=5$ (ακτίνα του δεύτερου), οπότε και οι τρεις πλευρές είναι $5$ cm.
\sol{15} Παίρνουμε τρία σημεία $\mathrm{Α}$, $\mathrm{Β}$, $\mathrm{Γ}$ της άκρης και σχεδιάζουμε τις χορδές $\mathrm{ΑΒ}$ και $\mathrm{ΒΓ}$. Φέρνουμε τις μεσοκάθετες των δύο χορδών (με διαβήτη). Το κέντρο ισαπέχει από τα άκρα κάθε χορδής, άρα βρίσκεται και στις δύο μεσοκαθέτους, δηλαδή στο σημείο τομής τους.
\sol{16} Επιστρέφει μετά από $6$ βήματα. Κάθε βήμα είναι χορδή ίση με την ακτίνα και μαζί με τις δύο ακτίνες σχηματίζει ισόπλευρο τρίγωνο, άρα η επίκεντρη γωνία είναι $60^\circ$. Έξι τέτοιες γωνίες κάνουν $360^\circ$.
\sol{17} Ελεύθερη εργασία. Για χάρτινο κύκλο, αν τον διπλώσεις σε δύο ίσα μέρη δύο φορές με διαφορετικές διπλώσεις, οι δύο γραμμές διπλώματος (διάμετροι) τέμνονται στο κέντρο.
\sol{18} Ελεύθερη κατασκευή. Η ακτίνα είναι $3{,}5$ cm. Με τον διαβήτη ανοιχτό $3{,}5$ cm βρίσκεις $6$ σημεία πάνω στον κύκλο και φτιάχνεις τόξα γύρω από το καθένα για τα πέταλα.
